\documentclass[11pt, a4paper]{article}

\usepackage{amsmath, amsthm, amssymb, amsfonts}
\usepackage{geometry}
\usepackage{graphicx}
\usepackage{hyperref}
\usepackage{natbib}
\usepackage{bm}
\usepackage{xcolor}
\usepackage{listings}
\usepackage{booktabs}

\geometry{margin=1in}

\newcommand{\pderiv}[2]{\frac{\partial #1}{\partial #2}}
\newcommand{\curl}{\nabla \times}
\renewcommand{\div}{\nabla \cdot}

\lstdefinestyle{code}{
    basicstyle=\ttfamily\small,
    breaklines=true,
    columns=fullflexible,
    frame=single,
    backgroundcolor=\color{gray!8}
}

\title{\textbf{HopfNet-MHD: Progress} \\
\large Milestones 1--4 --- Spectral Infrastructure, Gauge Projection, \\
Hopf Link Initialization, and the Nonlinear RHS Pipeline}
\author{Mohana Rangan Desigan}
\date{\today}

\begin{document}
\maketitle

\begin{abstract}
This report documents the implementation and validation of the first four
milestones of the HopfNet-MHD Phase 1 engine: a $128^3$ pseudo-spectral
incompressible Hall-magnetohydrodynamics (Hall-MHD) solver targeting Apple
Silicon (M4, 16\,GB unified memory). We report on (1) the spectral grid,
wavevector arrays, real-to-complex FFT wrappers, and Orszag $2/3$
dealiasing mask; (2) the Coulomb-gauge orthogonal projection tensor
$\mathcal{P}(\mathbf{k})$ and the spectral curl operator; (3) the
regularized Biot--Savart construction of the magnetic Hopf link initial
condition; and (4) the fully dealiased, gauge-consistent nonlinear
right-hand-side (RHS) pipeline for both the induction equation and an
incompressible Navier--Stokes velocity field, coupled through the Lorentz
force $\mathbf{J}\times\mathbf{B}$. Each milestone is accompanied by a
battery of unit tests verifying the relevant analytic identities to
machine precision on an $N=8$ test grid. All derivations reference the
governing equations established in the companion document
\emph{Mathematical Foundations of the HopfNet-MHD Engine}
\citep{foundation2026}.
\end{abstract}

\tableofcontents
\newpage

\section{Architecture Overview}
\label{sec:architecture}

The engine follows a hybrid C++/Python architecture, chosen to balance
numerical performance against development velocity on a single
16\,GB Apple Silicon workstation.

\begin{itemize}
    \item \textbf{C++ core} (\texttt{src/}, compiled via CMake +
    pybind11 into the extension module \texttt{hopfnet\_cpp}): owns the
    performance- and precision-critical primitives --- the spectral grid,
    the FFTW3-backed 3D real-to-complex transform pair, the Coulomb gauge
    projection tensor, the spectral curl operator, and the regularized
    Biot--Savart Hopf link initializer.
    \item \textbf{Python orchestration layer}
    (\texttt{python/hopfnet/rhs.py}): assembles the nonlinear RHS for both
    evolved fields ($\mathbf{A}$ and $\mathbf{v}$) by composing the C++
    primitives. Every floating-point-heavy operation (FFTs, projections,
    curls) executes inside compiled code; the Python layer performs only
    array bookkeeping via NumPy, which itself dispatches to C.
    \item \textbf{Test suite} (\texttt{tests/}): one file per milestone,
    run with \texttt{pytest}, executed on a small $N=8$ grid where
    analytic ground truth is tractable by hand.
\end{itemize}

This split keeps the stiff time-integration loop (Milestone 5--6) free to
be implemented entirely in C++ once the underlying operators are validated,
while allowing the nonlinear-term bookkeeping to be developed and debugged
interactively in Python first.

\section{Milestone 1: Spectral Grid and FFT Infrastructure}
\label{sec:milestone1}

\subsection{Mathematical Basis}

The pseudo-spectral method represents all fields as truncated Fourier
series (Sec.~2.1 of \citealp{foundation2026}, Eq.~10),
\begin{equation}
    \mathbf{A}(\mathbf{x}, t) = \sum_{\mathbf{k}} \hat{\mathbf{A}}(\mathbf{k}, t)
    e^{i \mathbf{k} \cdot \mathbf{x}},
\end{equation}
on a triply periodic domain of side $L = 2\pi$. For a real-valued field
sampled on $N^3$ points, the discrete wavevector components take the
standard FFT ordering $n \in \{0, 1, \dots, N/2, -(N/2-1), \dots, -1\}$,
with $k_i = (2\pi/L)\, n_i$.

\subsection{Implementation}

\texttt{src/spectral\_grid.hpp} implements the class \texttt{SpectralGrid},
which precomputes and stores:
\begin{itemize}
    \item 1D wavenumber arrays $k_x, k_y \in \mathbb{R}^N$ and
    $k_z \in \mathbb{R}^{N/2+1}$ (the last axis is truncated under the
    real-to-complex, or rFFT, layout, exploiting the Hermitian symmetry
    $\hat{\mathbf{A}}(-\mathbf{k}) = \hat{\mathbf{A}}^*(\mathbf{k})$ of a
    real field --- Sec.~2.4/\S\ref{sec:milestone2} of
    \citealp{foundation2026});
    \item the flattened array $|\mathbf{k}|^2 \in \mathbb{R}^{N \times N
    \times (N/2+1)}$;
    \item the Orszag $2/3$ dealiasing mask: a binary array that is $1$
    where $|n_x|, |n_y|, |n_z| \le \lfloor (2/3)(N/2) \rfloor$ and $0$
    otherwise.
\end{itemize}

\texttt{src/fft3d.hpp} implements \texttt{FFT3D}, a thin RAII wrapper
around a single pair of FFTW3 \texttt{r2c}/\texttt{c2r} plans created with
\texttt{FFTW\_MEASURE}. The forward transform returns the unnormalized
FFTW output; the inverse transform divides by $N^3$, so that
\texttt{inverse(forward(f))} $= f$ to machine precision. Both methods are
exposed to Python via pybind11 with NumPy array marshalling.

\subsection{Validation}

Six unit tests in \texttt{tests/test\_milestone1.py} establish:
\begin{enumerate}
    \item \textbf{Round-trip accuracy}: for a random field and for a
    constant field, \texttt{inverse(forward(f))} reproduces $f$ to
    $10^{-12}$.
    \item \textbf{Normalization}: a constant field $f \equiv 3.7$
    transforms to a single nonzero coefficient at $\mathbf{k}=0$ equal to
    $3.7 N^3$, confirming the $1/N^3$ inverse normalization.
    \item \textbf{Single-mode spectral purity}: $\cos(x)$ (integer
    wavenumber $n_x = \pm 1$) places $>99\%$ of its spectral energy at the
    expected indices $1$ and $N-1$.
    \item \textbf{Wavevector correctness}: for $N=8$, the computed $k_x,
    k_y, k_z$ arrays match the hand-derived FFT ordering
    $\{0,1,2,3,4,-3,-2,-1\} \times (2\pi/L)$ and $\{0,1,2,3,4\}\times
    (2\pi/L)$ respectively.
    \item \textbf{$|\mathbf{k}|^2$ consistency}: spot-checked against
    direct reconstruction from $k_x, k_y, k_z$.
    \item \textbf{Dealias mask correctness}: for $N=8$ the cutoff is
    $\lfloor (2/3)\cdot 4\rfloor = 2$; modes with $|n|\le 2$ are retained,
    $|n|\ge 3$ (including the Nyquist mode) are removed, and the mask is
    symmetric under $n \to -n$.
\end{enumerate}
All six tests pass.

\section{Milestone 2: Coulomb Gauge Projection and Spectral Curl}
\label{sec:milestone2}

\subsection{Mathematical Basis}

As derived in Sec.~2.3 of \citealp{foundation2026}, enforcing
$\div\mathbf{A}=0$ for all time reduces, in Fourier space, to applying the
orthogonal projection tensor
\begin{equation}
    \mathcal{P}_{ij}(\mathbf{k}) = \delta_{ij} - \frac{k_i k_j}{|\mathbf{k}|^2},
    \qquad
    \pderiv{\hat{\mathbf{A}}}{t} = \mathcal{P}(\mathbf{k})\,\hat{\mathbf{R}}.
    \label{eq:projection}
\end{equation}
At $\mathbf{k}=0$ this expression is $0/0$; following Sec.~2.3's
resolution, $\mathcal{P}(\mathbf{0}) \equiv \mathbb{I}$, since the
spatially uniform mean field trivially satisfies the gauge condition.

The spectral curl follows from $\nabla \to i\mathbf{k}$ in Fourier space:
\begin{equation}
    \hat{\mathbf{B}} = i\,\mathbf{k} \times \hat{\mathbf{A}},
    \qquad
    B_x = i(k_y A_z - k_z A_y),\quad
    B_y = i(k_z A_x - k_x A_z),\quad
    B_z = i(k_x A_y - k_y A_x).
    \label{eq:curl}
\end{equation}

\subsection{Implementation}

\texttt{src/projection.hpp} implements:
\begin{itemize}
    \item \texttt{project\_field(grid, Ax, Ay, Az)}: applies
    Eq.~\eqref{eq:projection} in place, with an explicit early-exit at
    $\mathbf{k}=0$ (identified by $|\mathbf{k}|^2 < 10^{-14}$) implementing
    $\mathcal{P}(\mathbf{0})=\mathbb{I}$.
    \item \texttt{spectral\_curl(grid, Ax, Ay, Az, Bx, By, Bz)}: implements
    Eq.~\eqref{eq:curl} component-wise, using the identity
    $i(a+ib) = -b + ia$ for the multiplication by $i$.
\end{itemize}
Both operate on \texttt{std::vector<std::complex<double>>} buffers of
length $N \times N \times (N/2+1)$ and are exposed to Python returning
NumPy \texttt{complex128} arrays of shape $(N,N,N/2+1)$.

\subsection{Validation}

Five unit tests in \texttt{tests/test\_milestone2.py}:
\begin{enumerate}
    \item \textbf{Gauge enforcement}: for a random complex vector field
    $\hat{\mathbf{R}}$, $\mathbf{k}\cdot\mathcal{P}(\mathbf{k})\hat{\mathbf{R}}
    = 0$ to $10^{-12}$ for every $\mathbf{k}\ne 0$.
    \item \textbf{Idempotence}: $\mathcal{P}(\mathcal{P}(\hat{\mathbf{R}}))
    = \mathcal{P}(\hat{\mathbf{R}})$, confirming $\mathcal{P}^2=\mathcal{P}$
    as required of a projection operator.
    \item \textbf{$\mathbf{k}=0$ identity}: the $\mathbf{k}=0$ mode of
    $\hat{\mathbf{R}}$ is passed through unchanged.
    \item \textbf{Curl correctness}: for the analytic field
    $\mathbf{A}=(0,0,\sin x)$, the spectral curl (via forward FFT $\to$
    \texttt{spectral\_curl} $\to$ inverse FFT) reproduces
    $\mathbf{B}=(0,-\cos x,0)$ to $10^{-12}$.
    \item \textbf{Automatic solenoidality}: for a random field, the curl
    output satisfies $\mathbf{k}\cdot\hat{\mathbf{B}}=0$ to $10^{-10}$,
    confirming $\div(\curl\mathbf{A})\equiv 0$ holds at the discrete level.
\end{enumerate}
All five tests pass.

\section{Milestone 3: Magnetic Hopf Link Initial Condition}
\label{sec:milestone3}

\subsection{Mathematical Basis}

The initial vector potential is constructed from two regularized
Biot--Savart current loops (Sec.~3 of \citealp{foundation2026}): Ring 1 in
the $xy$-plane centered at $(0,d,0)$, Ring 2 in the $yz$-plane centered at
$(0,-d,0)$, both of radius $R$. The infinite-wire singularity in the
Biot--Savart kernel is regularized by a finite core radius $a$
(Eq.~21),
\begin{equation}
    \mathbf{A}_{\mathrm{reg}}(\mathbf{r}) = \frac{\mu_0 I}{4\pi}
    \oint_C \frac{d\mathbf{l}'}{\sqrt{|\mathbf{r}-\mathbf{r}'|^2 + a^2}},
\end{equation}
giving the explicit component integrals of Eqs.~26--28, each a smooth
$2\pi$-periodic integral over the loop parameter.

\subsection{Implementation}

\texttt{src/hopf\_link.hpp} implements \texttt{hopf\_link::compute}, which:
\begin{itemize}
    \item evaluates each of the three component integrals (Eqs.~26--28)
    at every grid point using composite Simpson's rule on $[0,2\pi]$ with
    an even number of subintervals \texttt{n\_quad};
    \item places the grid on the centered domain
    $x_i = -L/2 + i\,(L/N)$, so the two interlocking rings sit
    symmetrically about the box center;
    \item returns three real $N\times N\times N$ arrays
    $(A_x, A_y, A_z)$.
\end{itemize}
The Coulomb-gauge projection (Milestone 2) is applied \emph{after}
construction, by transforming to Fourier space, calling
\texttt{project\_field}, and (if needed) transforming back ---
exactly as prescribed in Sec.~3.4 of \citealp{foundation2026}.

\subsection{Validation}

Four unit tests in \texttt{tests/test\_milestone3.py} (run at $N=8$,
$R=1$, $d=0.3$, $a=0.2$):
\begin{enumerate}
    \item \textbf{Finiteness and non-triviality}: all three components are
    finite everywhere (no NaN/Inf from the regularized $1/D$ kernels) and
    are not identically zero.
    \item \textbf{Quadrature convergence}: the maximum pointwise change in
    $A_x$ when doubling \texttt{n\_quad} from $64$ to $128$ is more than a
    factor of $2$ smaller than the change from $32$ to $64$ (Richardson-type
    convergence), and the $64\to128$ change is below $10^{-3}$ relative to
    the field's own magnitude --- consistent with Simpson's rule on a
    smooth, regularized integrand.
    \item \textbf{Coulomb gauge after projection}: after FFT $\to$
    \texttt{project\_field}, $\mathbf{k}\cdot\hat{\mathbf{A}} = 0$ to
    $10^{-12}$ for all $\mathbf{k}\ne 0$, exactly as required for the
    initial state to satisfy $\div\mathbf{A}=0$ to machine precision.
    \item \textbf{Physical content}: the resulting field
    $\mathbf{B}=\curl\mathbf{A}$ carries nonzero magnetic energy
    $\sum |\hat{\mathbf{B}}|^2 > 0$ (the Hopf link is not a trivial/null
    configuration) and is itself solenoidal,
    $\mathbf{k}\cdot\hat{\mathbf{B}}=0$ to $10^{-10}$.
\end{enumerate}
All four tests pass.

\subsection{Note on the $128^3$ Production Run}

The quadrature is $\mathcal{O}(N^3 \cdot n_{\mathrm{quad}})$. At $N=128$
with $n_{\mathrm{quad}}=64$, this is approximately $1.3\times 10^8$
evaluations of the regularized kernel per Cartesian component --- a
one-time initialization cost expected to take on the order of tens of
seconds on the M4, well within the 16\,GB memory budget
($\sim 260$\,MB per field).

\section{Milestone 4: The Nonlinear RHS Pipeline}
\label{sec:milestone4}

This milestone implements the right-hand side of both evolved equations
prior to the ETD split (Milestone 5). Two physical fields are evolved:
the vector potential $\mathbf{A}$ (Hall-MHD induction, Eq.~8) and an
incompressible velocity field $\mathbf{v}$ (momentum balance), coupled
through the Lorentz force $\mathbf{J}\times\mathbf{B}$.

\subsection{Mathematical Basis}

\paragraph{Induction RHS.} From Eq.~12 of \citealp{foundation2026}, the
raw forcing for $\mathbf{A}$ is
\begin{equation}
    \mathbf{R}_A = \mathbf{v}\times\mathbf{B}
    - \frac{d_i}{\rho}(\mathbf{J}\times\mathbf{B}) - \eta\mathbf{J},
    \label{eq:RA}
\end{equation}
which, after dealiasing, is projected via Eq.~\eqref{eq:projection} to
remove its longitudinal component, preserving $\div\mathbf{A}=0$.

\paragraph{Momentum RHS.} For an incompressible velocity field, the
nonlinear momentum forcing (advection plus Lorentz force, with the
pressure gradient eliminated by the same projection mechanism that
enforces $\div\mathbf{v}=0$) is
\begin{equation}
    \mathbf{R}_v = -(\mathbf{v}\cdot\nabla)\mathbf{v}
    + \frac{1}{\rho}(\mathbf{J}\times\mathbf{B}).
    \label{eq:Rv}
\end{equation}
The structural identity between Eqs.~\eqref{eq:RA} and \eqref{eq:Rv} ---
both are projected by the \emph{same} tensor $\mathcal{P}(\mathbf{k})$,
once enforcing the Coulomb gauge and once enforcing incompressibility ---
is a direct consequence of the fact that both constraints are statements
that a vector field's Fourier coefficients be transverse to $\mathbf{k}$
(Sec.~2.1 of \citealp{foundation2026}).

\paragraph{Dealiasing.} Each nonlinear product (the cross products
$\mathbf{v}\times\mathbf{B}$, $\mathbf{J}\times\mathbf{B}$, and the
advection term $(\mathbf{v}\cdot\nabla)\mathbf{v}$) is computed in real
space, transformed to Fourier space, and then masked by the Orszag $2/3$
dealiasing mask (Milestone 1) before being projected. This follows
Sec.~2's discussion: a quadratic nonlinearity convolves modes $k_1, k_2$
into $k_1+k_2$, which can exceed $k_{\max}$ and alias back into the
resolved spectrum unless removed.

\subsection{Implementation}

\texttt{python/hopfnet/rhs.py} implements:
\begin{itemize}
    \item \texttt{spectral\_derivative(grid, field\_hat, axis)}:
    $\partial_{x_{\mathrm{axis}}} \to i k_{\mathrm{axis}}$, used for the
    nine partial derivatives $\partial v_i/\partial x_j$ in the advection
    term;
    \item \texttt{compute\_B\_and\_J(grid, A\_hat)}: two applications of
    \texttt{spectral\_curl} (Milestone 2) give
    $\hat{\mathbf{B}}=\curl\hat{\mathbf{A}}$ and
    $\hat{\mathbf{J}}=\curl\hat{\mathbf{B}}$;
    \item \texttt{compute\_advection(grid, fft, v\_hat)}: computes
    $(\mathbf{v}\cdot\nabla)\mathbf{v}$ by transforming $\mathbf{v}$ and
    each $\partial v_i/\partial x_j$ to real space, forming the
    real-space sum $\sum_j v_j\,\partial v_i/\partial x_j$, transforming
    back, and applying the dealias mask;
    \item \texttt{cross\_real(a, b)}: elementwise real-space cross
    product, used for both $\mathbf{v}\times\mathbf{B}$ and
    $\mathbf{J}\times\mathbf{B}$;
    \item \texttt{compute\_R\_A} and \texttt{compute\_R\_v}: assemble
    Eqs.~\eqref{eq:RA} and \eqref{eq:Rv} respectively, each dealiasing
    every nonlinear product before the final call to
    \texttt{project\_field}.
\end{itemize}

This module is implemented in Python, but every floating-point-intensive
step --- the FFTs, the curls, and the projection --- is a call into the
compiled \texttt{hopfnet\_cpp} extension; the Python layer performs only
NumPy array arithmetic on already-transformed buffers.

\subsection{Validation}

Six unit tests in \texttt{tests/test\_milestone4.py}:
\begin{enumerate}
    \item \textbf{Spectral derivative}: $\partial_x \cos(x) = -\sin(x)$
    reproduced to $10^{-12}$ via \texttt{spectral\_derivative} followed by
    inverse FFT.
    \item \textbf{Dealiasing mechanism, demonstrated end-to-end}: the
    product $\cos(3x)\cos(2x) = \tfrac{1}{2}[\cos(5x)+\cos(x)]$ on an
    $N=8$ grid generates an aliased mode at integer wavenumber $n=-3$
    (from $5 \to 5-8$), which lies outside the $2/3$ cutoff
    ($\lfloor(2/3)\cdot 4\rfloor=2$). Before masking this mode is
    populated ($|\hat{f}|>10^{-6}$); after masking it is \emph{exactly}
    zero, while the surviving mode ($n=1$) is untouched.
    \item \textbf{Advection}: for $\mathbf{v}=(\sin x, 0, 0)$,
    $(\mathbf{v}\cdot\nabla)\mathbf{v} = (\tfrac{1}{2}\sin 2x, 0, 0)$,
    reproduced to $10^{-12}$.
    \item \textbf{Lorentz force chain}: for $\mathbf{A}=(0,0,\sin x)$
    (giving $\mathbf{B}=(0,-\cos x,0)$ and, by a second curl,
    $\mathbf{J}=(0,0,\sin x)$), the cross product
    $\mathbf{J}\times\mathbf{B} = (\tfrac{1}{2}\sin 2x, 0, 0)$ is
    reproduced to $10^{-12}$, validating the full curl
    $\to$ curl $\to$ cross-product chain.
    \item \textbf{$\mathbf{R}_A$ gauge consistency}: for random
    divergence-free $\hat{\mathbf{A}}, \hat{\mathbf{v}}$,
    $\mathbf{k}\cdot\hat{\mathbf{R}}_A = 0$ to $10^{-10}$.
    \item \textbf{$\mathbf{R}_v$ incompressibility consistency}: likewise,
    $\mathbf{k}\cdot\hat{\mathbf{R}}_v = 0$ to $10^{-10}$.
\end{enumerate}
All six tests pass.

\section{Summary of Test Coverage}

\begin{table}[h]
\centering
\begin{tabular}{@{}llc@{}}
\toprule
Milestone & Component & Tests \\
\midrule
1 & Spectral grid, FFT, dealias mask & 6 \\
2 & Coulomb gauge projection, spectral curl & 5 \\
3 & Hopf link Biot--Savart initialization & 4 \\
4 & Nonlinear RHS pipeline ($\mathbf{R}_A$, $\mathbf{R}_v$) & 6 \\
\midrule
\textbf{Total} & & \textbf{21} \\
\bottomrule
\end{tabular}
\end{table}

All 21 tests pass on $N=8$ test grids, exercising every operator that will
appear in the $128^3$ production loop: FFT round-trip and normalization,
the Coulomb-gauge projection tensor and its idempotence and $\mathbf{k}=0$
behavior, the spectral curl and its automatic solenoidality, the
regularized Biot--Savart initial condition and its convergence and gauge
consistency, and the fully dealiased, doubly-projected nonlinear RHS for
both the magnetic and velocity fields.

\section{Current Status and Next Steps}
\label{sec:status}

\subsection{What Is Validated}
\begin{itemize}
    \item All spectral differential operators ($\nabla$, $\curl$,
    $\div$ via projection) are implemated correctly to machine precision.
    \item The Coulomb-gauge projection $\mathcal{P}(\mathbf{k})$ correctly
    enforces $\div\mathbf{A}=0$ and, by the same construction,
    $\div\mathbf{v}=0$ for the incompressible velocity field --- including
    the $\mathbf{k}=0$ edge case.
    \item The Hopf link initial condition is generated, regularized, and
    gauge-projected, and carries genuine (non-trivial, solenoidal)
    magnetic energy.
    \item The full nonlinear forcing for both evolved fields,
    $\mathbf{R}_A$ and $\mathbf{R}_v$, is implemented with Orszag $2/3$
    dealiasing on every quadratic product, and both satisfy their
    respective transversality constraints after projection.
\end{itemize}

\subsection{What Remains for Phase 1}
\begin{itemize}
    \item \textbf{Milestone 5}: precompute the Kassam--Trefethen ETDRK4
    contour-integral coefficients $f_1, f_2, f_3, e_1=e^{\mathcal{L}\Delta
    t/2}, e_2=e^{\mathcal{L}\Delta t}$ for \emph{two} independent linear
    operators --- $\mathcal{L}_A = -\eta k^2 - \eta_4 k^4$ for the
    induction equation and $\mathcal{L}_v = -\nu k^2 - \nu_4 k^4$ for the
    momentum equation --- using the $M=32$, $R=1/\Delta t$ contour
    described in Sec.~4.3--4.4 of \citealp{foundation2026}.
    \item \textbf{Milestone 6}: assemble the ETDRK4 time-stepping loop
    coupling $\mathbf{R}_A$ and $\mathbf{R}_v$ from Milestone 4 with the
    two coefficient sets from Milestone 5, including checkpointing and
    online diagnostics ($\div\mathbf{A}$, $\div\mathbf{v}$, energy,
    $H_{\mathrm{gen}}$).
    \item \textbf{Milestone 7}: end-to-end short integration from the
    Hopf link IC at small grid size, as a fast regression smoke test.
    \item \textbf{$128^3$ production run}: once Milestones 5--7 pass,
    execute the full-resolution Hopf link evolution and verify the
    Algorithmic Validation Strategy of Sec.~6 of
    \citealp{foundation2026} (divergence-free to $10^{-13}$,
    $\mathcal{O}(\Delta t^4)$ temporal convergence, $H_{\mathrm{gen}}$
    conservation under $\eta_4=d_i=0$, and a clean $k^{-5/3}$ inertial
    range with no pile-up at $k_{\max}$).
\end{itemize}

\newpage
\begin{thebibliography}{99}

\bibitem[Desigan(2026)]{foundation2026}
Desigan, M. R. (2026). \textit{Mathematical Foundations of the HopfNet-MHD
Engine}. Phase 1.0 Documentation.

\bibitem[Kassam \& Trefethen(2005)]{kassam2005}
Kassam, A. K., \& Trefethen, L. N. (2005). Fourth-order time-stepping for
stiff PDEs. \textit{SIAM Journal on Scientific Computing}, 26(4),
1214-1233.

\bibitem[Abdelhamid et al.(2015)]{abdelhamid2015}
Abdelhamid, H. M., Kawazura, Y., \& Yoshida, Z. (2015). Hamiltonian
formalism of extended magnetohydrodynamics. \textit{Journal of Physics A:
Mathematical and Theoretical}, 48(23), 235502.

\bibitem[Orszag(1971)]{orszag1971}
Orszag, S. A. (1971). On the elimination of aliasing in finite-difference
schemes by filtering high-wavenumber components. \textit{Journal of the
Atmospheric Sciences}, 28(6), 1074.

\end{thebibliography}

\end{document}
